\title{LongRoPE: Extending LLM Context Window Beyond 2 Million Tokens}

\begin{abstract}
A substantial context window is an advantageous capability for large language models (LLMs). Nevertheless, limitations such as significant expenses associated with fine-tuning, limited availability of lengthy text corpora, and problematic values arising from unfamiliar token positions have restricted the expansion of context windows to approximately 128k tokens.
\end{abstract}

This paper presents LongRoPE, a method that, for the first time, scales the context window of pre-trained LLMs to an impressive \textbf{2048k} tokens. Notably, this is accomplished using no more than 1k fine-tuning steps at a maximum sequence length of 256k, all while preserving the model’s performance on its initial short context window. The approach is built upon three principal innovations: (i) we systematically uncover and utilize two types of non-uniformities in positional interpolation via an efficient search mechanism, yielding improved fine-tuning initialization and supporting an 8$\times$ context length increase without fine-tuning; (ii) we propose a staged extension framework that first fine-tunes the LLM to 256k tokens and subsequently applies a secondary positional interpolation on this model, thereby enabling a context window of 2048k; and (iii) we further adapt LongRoPE at an 8k length, restoring its effectiveness on short contexts. Comprehensive experiments conducted with LLaMA2 and Mistral over diverse tasks confirm the robustness and efficacy of our strategy. LongRoPE-based model extensions keep the original neural architecture intact with only minimal adjustments to the positional embedding and remain compatible with most existing optimizations. The source code will be accessible at \texttt{https://github.com/microsoft/LongRoPE}

\section{Introduction}

Although Large Language Models (LLMs) have achieved notable advancements across a range of tasks~\cite{gpt4,llama2}, their effectiveness is constrained by a limited context window—for instance, LLaMA2 is restricted to a maximum of 4096 tokens~\cite{llama2}. When the input exceeds this window, the model’s capabilities diminish because it encounters sequence positions outside those seen during training. This limitation becomes particularly problematic in critical applications such as in-context learning with many examples~\cite{Huang2023FewerIM} and the operation of LLM agents~\cite{park2023generative,madaan2023selfrefine}.

Recent studies have demonstrated that the context window of a pre-trained LLM can be enlarged to approximately 128k through fine-tuning with extended textual inputs~\cite{longlora,pi,yarn,zhang2024soaring,liu2023scaling}. Despite these advances, further increasing the context window encounters several significant challenges. First, the introduction of numerous uninitialized position indices creates catastrophic values, resulting in distributional shifts that hamper the convergence of fine-tuning~\cite{pi}. This complication becomes pronounced when expanding from 4k to over 1000k, since more than 90\% of the positions are entirely new. Second, effective fine-tuning typically requires access to text samples of equivalent lengths; unfortunately, the availability of such lengthy texts—particularly those greater than 1000k—in existing datasets is scarce. Furthermore, training with extremely long textual data imposes heavy computational demands, necessitating substantial amounts of training time and GPU resources. Third, scaling to ultra-long context windows causes the model's attention to be distributed over a vast range of token positions, ultimately reducing performance on tasks requiring shorter contexts~\cite{pi}.

To address the initial challenge, a common strategy involves interpolating RoPE positional embeddings~\cite{rope,pi}, which rescales incoming position indices to align with the pretrained domain, as illustrated in Fig.~\ref{fig:overview}. The Position Interpolation (PI) method~\cite{pi} achieves this by linearly adjusting RoPE’s rotary angles according to the expansion factor. The NTK approach~\cite{ntk,dynamicntk} employs non-uniform interpolation and extrapolation across different RoPE dimensions. YaRN~\cite{yarn} divides RoPE dimensions into three distinct frequency-based categories, utilizing extrapolation, NTK interpolation, and linear interpolation in each group. Despite these techniques, positional embedding within the Transformer framework embodies a nuanced, non-uniform distribution of information entropy. Existing methods fail to exploit this intricate non-uniformity fully, resulting in information degradation and restricting the attainable context window.

Section~\ref{sec:analysis} presents two principal empirical insights: \textbf{(1)} Robust positional interpolation necessitates addressing two distinct types of non-uniformity—specifically, differing RoPE dimension ranges as well as token locations within the sequence. Notably, positions with lower RoPE dimensions and early token indices generally require reduced degrees of interpolation; however, the precise approach must be adapted to the intended maximum sequence length. \textbf{(2)} Integrating these identified non-uniformities into the design of positional interpolation procedures makes it possible to better preserve critical information from the original RoPE encoding, especially in dimensions and positions of heightened significance. This strategy mitigates the typical degradation stemming from positional interpolation and consequently offers a more effective starting point for subsequent fine-tuning. In addition, it facilitates up to an 8$\times$ extension in input length, even in contexts where additional fine-tuning is not performed.

Driven by these observations, we present \textbf{LongRoPE}, a robust approach capable of pushing the LLM context window past 2 \emph{million} tokens. The design of LongRoPE leverages three principal advancements. First, LongRoPE takes advantage of multidimensional, non-uniform positional interpolation by selecting optimal rescale factors for RoPE's rotational angles within each RoPE dimension, based on individual token positions. Because the task of determining these rescale factors is exponentially more complex with increasing extension ratios, LongRoPE employs an evolutionary search mechanism augmented with two optimization strategies to significantly enhance the efficiency of the search. An illustration of the resultant rescaled RoPE is depicted in Fig.~\ref{fig:overview}.

Subsequently, {LongRoPE} adopts a resource-efficient and incremental extension method to realize a context window of 2048k, circumventing the need for fine-tuning on exceedingly long input sequences, which are rare and challenging to obtain. This approach initiates by identifying a suitable 256k length threshold using the original pre-trained LLM, followed by fine-tuning at this scale. Given that the non-uniform positional interpolation enables an eightfold expansion without additional fine-tuning, a further optimization is performed by determining new RoPE rescaling parameters on the already extended, fine-tuned LLM. Through this process, the framework attains a 2048k context window for both LLaMA2 and Mistral~\cite{mistral}.

Lastly, to prevent a decline in performance within the original, more limited context window, {LongRoPE} refines the RoPE rescale factors for the expanded LLM model. In line with the earlier process of scaling the context size up from 256k to 2048k, we apply a downscaling to 4k and 8k windows in the 256k fine-tuned LLM, utilizing our search algorithm to reduce the degree of positional interpolation. When the input sequence during inference falls below 8k in length, RoPE is updated using the rescale factors identified by this search.

Comprehensive evaluations spanning multiple LLMs and a range of long-context tasks highlight the strong performance of our approach. Our results indicate that LongRoPE consistently preserves low perplexity across evaluation sequences extending from 4k up to 2048k, achieves passkey retrieval accuracies exceeding 90\%, and matches the accuracy of baseline models on standard tasks constructed within a 4096 context window. LongRoPE is broadly compatible with any LLM utilizing RoPE embedding. Both our codebase and the LongRoPE-2048k models will be made publicly available.

\section{Non-uniformity in Positional Interpolation}
\label{sec:analysis}

\subsection{Preliminary}

Transformer models require explicit positional information, often in the form of position embedding, to represent the order of input tokens. Our work focuses on the RoPE~\cite{rope} position embedding, which is widely used in recent  LLMs. For a token at position index $n$, its corresponding RoPE encoding can be simplified as follows:

\begin{equation}
		\fontsize{7.8}{7.8} \selectfont
	\label{eq:rope}
	\left[ cos(n\theta_0), sin(n\theta_0), cos(n\theta_1), \cdot\cdot\cdot, cos(n\theta_{d/2-1}), sin(n\theta_{d/2-1}) \right]   
\end{equation}

Here, $d$ refers to the size of the embedding vector, $n\theta_i$ denotes the rotation angle applied to the token located at position $n$, and $\theta_i=\theta^{-2i/d}$ specifies the frequency used for rotation. By default, RoPE adopts a base $\theta$ value of 10000.

\textbf{\emph{Context window extension ratio $s$} and positional interpolation}. The parameter $s$ is introduced as the quotient of the extended context window length $L'$ and the original length $L$, expressed as $s=\frac{L'}{L}$.

To extend the context window from $L$ to $L'$, current positional interpolation methods suggest downscaling rotation frequencies $\theta_i$ based on extension ratio $s$. Let $\beta=\theta^{2/d}$, and $\lambda$ denote the actual rescale factor related to $s$, we   unify these positional interpolation methods as follows:

\begin{equation}	
	\fontsize{7.5}{7.5} \selectfont
	\label{eq:unified_rope}
	\begin{aligned}
		\left[cos\left(\frac{n}{\lambda(\beta)^0}\right), sin \left(\frac{n}{\lambda(\beta)^0}\right), cos\left(\frac{n}{\lambda(\beta)^1}\right), \cdot\cdot\cdot,  sin \left(\frac{n}{\lambda(\beta)^{d/2-1}}\right) \right]   
	\end{aligned}
\end{equation}

\textbf{Linear positional interpolation (PI)}. PI~\cite{pi} proposes using a linear scaling of positional indices inside the range for which the model was originally trained. Specifically, for an intended extension ratio $s$, all position rotation angles are uniformly scaled down by $\lambda = s$ across each RoPE dimension. This uniform compression results in the positional encodings becoming densely packed, which in turn impairs the model’s capacity to differentiate between tokens at nearby positions. As a consequence, PI generally exhibits diminished effectiveness, particularly when extending to large values of $s$.

\textbf{NTK-based interpolation and extrapolation}. Research by~\cite{ntk,dynamicntk} approaches RoPE from the angle of information encoding, leveraging Neural Tangent Kernel (NTK) theory~\cite{jacot2018neural,tancik2020fourier} for analysis. To address the problem of overlapping positions found in PI, they propose to balance interpolation demands across RoPE’s multiple dimensions. Specifically, their method applies reduced scaling to dimensions associated with higher frequencies and increased scaling to those linked with lower frequencies. This enables effective interpolation and extrapolation of positions, following the form $\lambda=s^{i}$. The enhanced dynamic NTK~\cite{dynamicntk} further introduces position-dependent adaptation of the extension ratio in line with the current length of the sequence. In contrast to PI, which requires additional fine-tuning, NTK-aware approaches permit the extension of context windows without the need for fine-tuning; however, such extensions are generally capped at a maximum ratio of 4$\times$.

\textbf{YaRN}~\cite{yarn} brings notable advancements to positional interpolation by segmenting RoPE dimensions into three groups based on frequency, each assigned a unique interpolation method. Specifically, dimensions with high frequencies are handled via extrapolation ($\lambda$=1), low-frequency dimensions adopt linear interpolation (PI), and those with intermediate frequencies are managed using NTK. The core innovation in YaRN stems from its partitioning of RoPE dimensions, a process currently guided by manual, empirically-driven decision-making. Consequently, this approach might yield less-than-ideal results when applied to newly developed LLMs.

\subsection{Study on Non-uniform Positional Interpolation}

Drawing on insights from NTK and YaRN, we observe that their improvements stem from introducing non-linearity, particularly through the treatment of frequency variations across RoPE dimensions to enable customized interpolation and extrapolation. Nonetheless, prevailing non-linear methods are largely based on handcrafted heuristics. This observation leads us to two fundamental questions: (1) Does the existing positional interpolation represent the best possible approach? (2) Might there be alternative, yet-unidentified forms of non-linearity worth exploring?

In order to investigate these issues, we employ evolution search (refer to Sec.~\ref{sec:method}) to identify more effective non-uniform positional interpolation schemes for LLaMA2-7B. The optimization process leverages perplexity as a metric and is conducted with 5 randomly selected samples from the PG19~\cite{pg19} validation set. Our empirical evaluation leads to several noteworthy insights.

\textbf{Finding 1}: \textit{RoPE dimensions exhibit substantial non-uniformities, which are not effectively handled by current positional interpolation methods.}

We optimize $\lambda$ for each RoPE dimension as specified in Eq.~\ref{eq:unified_rope}. In Table~\ref{tbl:exp1}, we present the perplexity results for LLaMA2-7B using various approaches on the PG19 and Proof-pile~\cite{proof-pile} evaluation sets, conducted without further fine-tuning. The proposed optimization yields clear perplexity reductions, indicating that traditional linear (PI) and alternative non-uniform interpolation schemes (Dynamic-NTK and YaRN) are not optimal. Interestingly, YaRN performs worse than PI and NTK on the PG19 dataset, as it does not successfully achieve the target context window size for models that have not been fine-tuned. For instance, with YaRN, the perplexity increases sharply past the 7k token mark within an 8k token context window.

In our investigation, we found that the rescaling factors $\lambda$ in Eq.~\ref{eq:unified_rope} are not uniform across dimensions, as opposed to the constant scale $s$ applied in PI, the formula used in NTK, or the group-wise approach of YaRN. The introduction of these variable scaling factors leads to notable enhancements in LLaMA2's language modeling accuracy (measured by perplexity) when handling context lengths of 8k and 16k, all without requiring additional fine-tuning. This improvement arises because the adapted positional embeddings better retain the characteristics of the original RoPE, with key dimensions especially well preserved, thereby helping the LLM more effectively differentiate among nearby token positions.

\textbf{Finding 2}: \textit{RoPE for the initial tokens in the input sequence should be extrapolated with less interpolation.} 

We postulate that the first $\hat{n}$ tokens within input sequences require diminished interpolation from RoPE. This stems from their typically higher attention scores, rendering them particularly influential within attention mechanisms, as evidenced by findings from Streaming LLM~\cite{streamingllm} and LM-Infinite~\cite{han2023lminfinite}. To evaluate this premise, we expand the context length to 8k and 16k by applying PI and NTK, while explicitly preserving the initial $\hat n$ tokens (ranging from 0, 2, ..., 256) without any interpolation. Setting $\hat n = 0$ effectively returns to the baseline PI and NTK. The results in Table~\ref{tbl:tokenposition} reveal two principal insights: \textbf{(1)} foregoing position interpolation for the leading tokens leads to notable gains in performance, and \textbf{(2)} the ideal quantity for $\hat n$ varies according to the length of the extended context.

\textbf{Finding 3}: \textit{Non-uniform positional interpolation effectively extends LLM context window in both fine-tuning and non-fine-tuning settings.} 

Our experiments demonstrated that implementing the discovered non-uniform position interpolation noticeably boosts extension performance at 8k and 16k context lengths even without fine-tuning. However, successfully scaling to longer contexts necessitates fine-tuning. Accordingly, we applied fine-tuning to LLaMA2-7B using the optimized RoPE settings for an extended 64k context window (refer to the Appendix for detailed configuration). As reported in Table~\ref{tbl:finetune}, our approach achieves markedly superior results relative to PI and YaRN, both prior to and following the fine-tuning process. This enhancement stems from leveraging non-uniform positional interpolation, which reduces information degradation and offers a superior starting point for fine-tuning.

\textbf{Summary}. This work identifies two non-uniform aspects in RoPE—its dimensional distribution and token placement. Effectively exploiting these non-uniformities within positional interpolation leads to substantial gains in context extension capabilities for LLMs.

\section{LongRoPE}

\label{sec:method}
Building on these insights, we propose LongRoPE. This method initially implements a fast search strategy designed to leverage the two observed non-uniformities, and subsequently applies it to expand the context window of LLMs to more than 2 million tokens.

\subsection{Problem Formulation}

These two forms of non-uniformity expand the solution space considerably and add layers of complexity to the optimization process. To tackle this, we recast the multidimensional non-uniform position interpolation optimization challenge as a search problem.

For a LLM targeting a  context window size of $L'$ and lengthy input documents $\mathbf{X}$, where each $\mathbf{x}\in\mathbf{X}$ surpasses $L'$ in token length, we denote the original rotary angle of the $i^{th}$ dimension in RoPE embedding at token position $n$ as  $\frac{n}{\beta^i}$. The optimization problem is then formulated as follows:

\begin{equation}
\begin{gathered}
		\label{eq:problem}

\underset {\mathbf{x}\in\mathbf{X}; \, |\mathbf{x}|\ge L'}{ \operatorname {arg\,min} } \,  \mathcal{L}\left( \text{LLM}(\text{RoPE}, \mathbf{X})\right), \text{where}

\\
\scalebox{0.93}{
$\underset {\begin{subarray}{l} i=0,\ldots,\frac{d}{2}-1;\\ n\in[0,|\mathbf{x}|);\end{subarray}}{\text{RoPE}_(n)} = \left[\ldots, \cos\left(\mathbb{I}(\hat{\lambda}_{i}, \hat{n}) \frac{n}{\beta^i}\right), \sin\left(\mathbb{I}(\hat{\lambda}_{i}, \hat{n}) \frac{n}{\beta^i}\right), \ldots \right] $}\\
	\text{where} \quad \mathbb{I}(\hat{\lambda}_{i}, \hat{n}) = \begin{cases}
	1 &\text{if} \quad n<\hat{n} \\
	{\frac{1}{\lambda}_i} &\text{if} \quad n\geq\hat{n}
	\end{cases}
	\end{gathered}
\end{equation}

In this approach, a set of scaling factors $\mathbb{I}(\hat{\lambda}_{i},\hat{n})$ are introduced to address both types of non-uniformities present. Here, $\hat{\lambda_i}$ refers to irregularities among RoPE dimensions, and $\hat{n}$ signifies non-uniformity with respect to token positions. The scaling function $\mathbb{I}(\hat{\lambda}_{i},\hat{n})$ is specifically used to adjust the rotation angle in the $i^{th}$ RoPE dimension, employing $\hat{\lambda}_{i}$ as the scaling coefficient and $\hat{n}$ as a positional threshold for tokens. For token positions less than $\hat{n}$, $\hat{\lambda}_{i}$ does not influence the computation and the standard RoPE rotary angle $\frac{n}{\beta^i}$ remains unchanged. However, once the token position $n$ is greater than or equal to $\hat{n}$, the scaling factor is then applied.

With a given target context window length of $L'$, our task is to determine the set of rescaling factors ($\mathbb{I}(\hat{\lambda}_{0},\hat{n})$, $\mathbb{I}(\hat{\lambda}_{1},\hat{n})$, ...$\mathbb{I}(\hat{\lambda}_{i},\hat{n})$...) for each dimension of the RoPE mechanism, spanning from the first to the $d^{th}$. By applying these rescale factors to the $\text{RoPE}$ within the target $\text{LLM}$, the aim is to reduce the next token prediction loss, denoted as $\mathcal{L}$ (that is, perplexity), on sample inputs $\mathbf{X}$ of token length $L'$.

\subsection{Searching the Non-uniform Position Interpolation}

To address the challenge presented in Eq.~\ref{eq:problem}, we propose an efficient and straightforward approach that identifies the most suitable RoPE rescaling coefficients. This method aims to maximize utilization of the complex, multidimensional irregularities inherent in position embeddings.

\textbf{Search space}. Our search space is constructed to encompass both types of non-uniformities, resulting in a substantially large configuration set. The details of this search space are presented in Table~\ref{tbl:searchspace}. Rather than seeking $\mathbb{I}(\hat{\lambda}_{i},\hat{n})$ directly, we simplify the design by searching over the individual rescaling factor $\lambda_i$ for each RoPE embedding dimension, and the parameter $\hat{n}$. Here, $\hat{\lambda}_{i}$ is defined as $1/\lambda_i$. According to Table~\ref{tbl:searchspace}, the range of $\lambda_i$ is set from 1.0 (representing direct extrapolation) up to $s\times1.25$ (allowing broader interpolation than PI), with increments of 0.01. In this context, $s$ denotes the ratio for extending the target context window.

The parameter $\hat{n}$ determines how many of the initial token positions are preserved using their original RoPE embedding, rather than undergoing position interpolation. In practice, $\hat{n}$ is chosen from the set \{0, 1, 2, 4, 8, 12, 16, 20, 24, 28, 32, 64, 128, 256\}. Setting $\hat{n}=0$ implies that rescale factors, as found through the search process, are applied to every token position.

\textbf{Evolution-based search}. The search space described in Table~\ref{tbl:searchspace} encompasses a vast array of positional interpolation methods, creating considerable obstacles for effective search strategies. As an illustration, extending by $s=4\times$ yields $400^{128/2}\times14$=4$\times10^{167}$ possible configurations. The search space grows exponentially with increasing extension ratios. To overcome these hurdles, we employ an evolution search strategy~\cite{guo2020single} supplemented with two optimization approaches that significantly enhance the efficiency of exploration. Algorithm~\ref{alg:search} outlines the comprehensive search workflow.

\textit{Optimized initial population generation}. Rather than assigning all $P$ rescale factors at random, our approach directly includes the RoPE rescale factors for PI, NTK, and YaRN among the first individuals in the population. For the other $P-3$ members, we introduce diversity by applying random mutations to these three rescale factors, each with probability $p$.

\textit{Monotonically non-decreasing constraint}. Once the initial population has been established, LLM perplexity is assessed for every member. To do this, the appropriate RoPE rescale factors are applied to the target LLM, followed by calculating the perplexity on the input $\mathbf{X}$. The individuals with the best (lowest) perplexity scores are selected as parents for the evolutionary process. Nevertheless, the vastness of the search space means that simple mutation and crossover operations may yield suboptimal candidates, resulting in unnecessary perplexity evaluations. This inefficiency is exacerbated when $L'$ is large, since each perplexity computation requires a significant amount of time for inference.

To mitigate this issue, we enforce a constraint that ensures the rescaled RoPE factors are monotonically non-decreasing: $\lambda_i \le \lambda_{i+1}$.
Only those RoPE configurations that adhere to this monotonicity are considered for perplexity testing in LLMs, thereby streamlining and accelerating the search process. In particular, we stipulate that $\lambda_i$ grows in a monotonic fashion across the RoPE dimensions ($i = 0, \ldots, 63$). This requirement is motivated by Neural Tangent Kernel (NTK) theory~\cite{jacot2018neural,tancik2020fourier,ntk}, which posits that dimensions corresponding to higher frequencies (lower index values) necessitate less interpolation—implying smaller $\lambda_i$ values—while dimensions associated with lower frequencies (higher index values) allow for increased interpolation, reflected in larger $\lambda_i$ values.

\begin{algorithm}[t]
	\small
	\caption{The search algorithm}
	\textbf{Input:} target LLM, input samples $\mathbf{X}$, population size $P$, mutation size $N_1$, crossover size $N_2$, maximum iteration count $\mathcal{T}$, mutation probability $p$\\

	\begin{algorithmic}[1]
		\label{alg:search}
		\STATE Top-k $\gets$ $\phi$;
		\STATE $\text{P}_0 \gets$ \textit{Initialize\_population\_with\_optimization} ($P$, $\mathbf{X}$, $p$);
		\FOR{$i = 1$ to $\mathcal{T}$}
		\STATE \textit{Compute\_perplexity} (LLM, $\text{P}_{i-1}$, $\mathbf{X}$);
		\STATE Top-k $\gets$ \textit{Update\_Topk} (Top-k, $\text{P}_{i-1}$);
		\STATE $\text{P}_{mutation} \gets$ \textit{Mutation\_with\_mono\_constraint} (Top-k, $N_1$, $p$);
		\STATE $\text{P}_{crossover} \gets$ \textit{Crossover\_with\_mono\_constraint} (Top-k, $N_2$);
		\STATE $\text{P}_i \gets \text{P}_{mutation} \cup \text{P}_{crossover} \cup$ Top-k;
		\ENDFOR
		\STATE Return the individual in Top-k that exhibits the minimum perplexity;
	\end{algorithmic}
\end{algorithm}

\textbf{8$\times$ extension without fine-tuning}. Through evolutionary search, our approach efficiently determines optimal non-uniform RoPE rescale factors, maintaining crucial dimensions and token positions to reduce information degradation due to interpolation. As illustrated in Fig.\ref{fig:non-ft}, this technique enables LLaMA2's context window to expand from 4k to 32k tokens without requiring any fine-tuning. By comparison, current methods—including PI, as well as non-uniform NTK and YaRN—lead to a significant increase in perplexity beyond a 2$\times$ window extension.

\subsection{Extending LLM Context Window to 2048K}
\label{sec:extendto2048k}

\textbf{Progressive extension to 2048k}. In this section, we present our approach for scaling the context window of pre-trained LLMs from the standard 4k up to over 2048k tokens. Our experiments show that non-uniform positional interpolation allows for an 8$\times$ context window increase without the need for fine-tuning. However, when aiming for much higher extension ratios, such as 512$\times$, fine-tuning becomes unavoidable. A common technique involves identifying appropriate RoPE rescaling parameters for the desired 2048k window and then performing fine-tuning. Nevertheless, this procedure is constrained by the significant computational costs involved. Furthermore, in our observations, effectively fine-tuning LLMs with such extensive context window extensions remains a difficult task (refer to Appendix).

Luckily, LongRoPE demonstrates efficacy in enhancing both the base and the fine-tuned extended LLM models. To this end, we propose a streamlined, stepwise approach that reaches the desired 2048k context length using merely 1k fine-tuning iterations, while maintaining a maximum training sequence length of 256k.

$\Diamond$ \textit{Utilizing LongRoPE search to broaden pre-trained LLM to 256k.} Using LLaMA2 as a case study, we explore searches targeting context window sizes of 128k and 256k. This phase involves extension ratios of 32$\times$ and 64$\times$ correspondingly.

$\Diamond$ \textit{Fine-tuning to 256k}. Subsequently, we adapt the pre-trained LLM to support a 256k context window. In detail, our approach involves initially fine-tuning LLaMA2 for 400 iterations employing RoPE rescaled factors set to 128k. Afterward, we update the checkpoint by substituting the RoPE rescaled factors with those corresponding to 256k, followed by 600 further fine-tuning steps. Compared to direct fine-tuning at 256k, this strategy demonstrates superior efficiency.

$\Diamond$ \textit{LongRoPE search enables expansion of a fine-tuned extended LLM to 2048k}. In the last stage, we apply an additional search procedure to the LLM that has been fine-tuned at 256k context length. This allows the model to reach an exceptionally wide context window of 2048k, achieved without additional rounds of fine-tuning. The process yields a total extension factor of $512\times$.

\textbf{Restoring performance in reduced context windows}.  Following the expansion to a substantially larger 2048k context window, there is an observable decrease in performance in the initial context range. This decline is a documented consequence of positional interpolation~\cite{pi}, which compresses the position embeddings for the original context into a relatively confined space, thereby impairing the model's language capabilities. When the context window is stretched by a factor of 512$\times$, the positional embeddings within the original 4k window are densely packed together.

To address this issue, an additional evolution search is conducted on the extended LLM, specifically to fine-tune the RoPE rescale factors for shorter context lengths (such as 4k and 8k). For these shorter lengths, we constrain the upper bound of the searched $\lambda$ since they demand less positional interpolation. At inference time, the LLM autonomously adapts the relevant RoPE rescale factors according to the current context length.

\section{LongRoPE}

\label{sec:method}
Building on these insights, LongRoPE is proposed. This method begins by implementing a streamlined search algorithm designed to leverage both forms of non-uniformity, subsequently enabling the extension of the LLM context window to accommodate more than 2 million tokens.

\subsection{Problem Formulation}

These two forms of non-uniformity greatly expand the solution space and complicate the optimization process. To mitigate these challenges, we conceptualize the optimization of multidimensional non-uniform position interpolation as a search problem.

For a LLM targeting a  context window size of $L'$ and lengthy input documents $\mathbf{X}$, where each $\mathbf{x}\in\mathbf{X}$ surpasses $L'$ in token length, we denote the original rotary angle of the $i^{th}$ dimension in RoPE embedding at token position $n$ as  $\frac{n}{\beta^i}$. The optimization problem is then formulated as follows:

\begin{equation}
\begin{gathered}
		\label{eq:problem}

\underset {\mathbf{x}\in\mathbf{X}; \, |\mathbf{x}|\ge L'}{ \operatorname {arg\,min} } \,  \mathcal{L}\, \Big( \text{LLM} (\text{RoPE}, \mathbf{X})\Big), \text{with} 
\\
\scalebox{0.93}{
$\underset {\begin{subarray}{l} i=0,\ldots, \frac{d}{2}-1; \\ n\in [0, |\mathbf{x}|); \end{subarray}}{\text{RoPE}_(n)}= \left[\ldots, \cos\left(\mathbb{I}(\hat{\lambda}_{i}, \hat{n}) \frac{n}{\beta^{i}}\right), \sin\left(\mathbb{I}(\hat{\lambda}_{i}, \hat{n}) \frac{n}{\beta^{i}}\right), \ldots\right]$}\\
	\text{with} \; \mathbb{I}(\hat{\lambda}_{i},\hat{n})=\begin{cases}
	1&\mbox{ if $n<\hat{n}$} \\
	\frac{1}{\lambda}_{i}&\mbox{ if $n\geq\hat{n}$}
	\end{cases}
	\end{gathered}
\end{equation}

A set of rescale factors, $\mathbb{I}(\hat{\lambda}_{i},\hat{n})$, is introduced to account for both types of non-uniformities. Here, $\hat{\lambda_i}$ corresponds to non-uniformity across RoPE dimensions, while $\hat{n}$ represents variations in token positions. In practice, the function $\mathbb{I}(\hat{\lambda}_{i},\hat{n})$ serves to adjust the rotation angle for the $i^{th}$ RoPE dimension, with $\hat{\lambda}_i$ acting as the scaling factor and $\hat{n}$ as a threshold for token positions. For tokens with positions less than $\hat{n}$, the scaling factor $\hat{\lambda}_{i}$ is not applied, and the original RoPE rotary angle $\frac{n}{\beta^i}$ is retained. However, for token positions $n\ge\hat{n}$, the scaling factor comes into effect.

Assuming a desired context window of length $L'$, the task involves determining the optimal set of rescale coefficients ($\mathbb{I}(\hat{\lambda}_{0},\hat{n})$, $\mathbb{I}(\hat{\lambda}_{1},\hat{n})$, ... , $\mathbb{I}(\hat{\lambda}_{i},\hat{n})$, ...) across RoPE dimensions from the $1^{st}$ to the $d^{th}$. By tuning these parameters, the resulting target $\text{LLM}$, incorporating the adjusted $\text{RoPE}$, aims to minimize next-token prediction error $\mathcal{L}$ (i.e., perplexity) when processing input sequences $\mathbf{X}$ of token length $L'$.

\subsection{Searching the Non-uniform Position Interpolation}

To address the challenge outlined in Eq.~\ref{eq:problem}, we present a straightforward and powerful approach that seeks the best RoPE rescale parameters, enabling comprehensive utilization of the diverse, multidimensional variations within position embedding.

\textbf{Search space}. We construct an extensive search space that incorporates both non-uniformities, as outlined in Table~\ref{tbl:searchspace}. In particular, our approach enables the exploration of tailored rescaling factors for each RoPE embedding dimension. To streamline the search space configuration, we choose to optimize $\lambda_i$ and $\hat{n}$, rather than $\mathbb{I}(\hat{\lambda}_{i},\hat{n})$, with $\hat{\lambda}_{i}$ defined as $1/\lambda_i$. According to Table~\ref{tbl:searchspace}, the search for $\lambda_i$ is conducted over a range starting from 1.0 (representing direct extrapolation) up to $s\times1.25$ (which exceeds the interpolation of PI), incremented by 0.01, where $s$ indicates the desired context window expansion ratio.

The parameter $\hat{n}$ determines how many token positions at the beginning are preserved with their original RoPE embeddings, bypassing position interpolation. In practice, we explore values for $\hat{n}$ from the set \{0, 1, 2, 4, 8, 12, 16, 20, 24, 28, 32, 64, 128, 256\}. If $\hat{n}=0$, rescale factors discovered through search are applied to every token position.

\textbf{Evolution-based search}. The search space outlined in Table~\ref{tbl:searchspace} encompasses a wide array of positional interpolation configurations, resulting in complex and large-scale combinatorial possibilities that are difficult to explore efficiently. For instance, increasing the extension to $s=4\times$ generates $400^{128/2}\times14$=4$\times10^{167}$ potential options. As the extension ratio grows, the size of the search space rises exponentially. To make exploration feasible, we apply evolution search~\cite{guo2020single} and incorporate two key optimization strategies that substantially enhance the search process. Algorithm~\ref{alg:search} provides an overview of this search methodology.

\textit{Optimized initial population generation}. Rather than starting with an entirely random selection of $P$ rescale factors for the population, we directly include the three specific RoPE rescale factors—PI, NTK, and YaRN—as distinct members of the initial set. The other $P-3$ individuals are then created by applying random mutations to these three rescale factors, each with probability $p$.

\textit{Monotonically non-decreasing constraint}. Once the initial population is established, we evaluate the LLM perplexity associated with each candidate. For this process, the appropriate RoPE rescale factors are applied to the target LLM, and the perplexity of the input $\mathbf{X}$ is then measured. The individuals with the lowest perplexity, specifically the top-k, are selected as parents for subsequent evolutionary steps. Nevertheless, the immense size of the search space means that unrestrained mutation and crossover operations may frequently yield suboptimal candidates, necessitating numerous unnecessary perplexity evaluations. This inefficiency is further exacerbated for large values of $L'$, since each perplexity assessment is computationally intensive.

To mitigate this issue, we enforce a constraint ensuring the RoPE rescaled factors $\lambda_i$ form a non-decreasing sequence, that is, $\lambda_i\le \lambda_{i+1}$. We limit our perplexity evaluations on LLMs to RoPE configurations adhering to this restriction, which leads to a substantial reduction in search complexity. Concretely, we specify that $\lambda_i$ must increase monotonically across the RoPE dimension indices ($i = 0, ..., 63$). This requirement is motivated by NTK theory~\cite{jacot2018neural,tancik2020fourier,ntk}, which argues that components in lower dimensions (associated with higher frequency) necessitate less interpolation (corresponding to a smaller value of $\lambda_i$), while those in higher dimensions (with lower frequency) permit greater interpolation (implying larger $\lambda_i$).

\begin{algorithm}[t]
	\small
	\caption{The search algorithm}
	\textbf{Input:} target LLM, input samples $\mathbf{X}$, population size $P$, mutation size $N_1$, crossover size $N_2$, max iterations $\mathcal{T}$, mutation probability $p$\\

	\begin{algorithmic}[1]
		\label{alg:search}
		\STATE Top-k $\gets$ $\phi$;	
		\STATE $\text{P}_0$ $\gets$ \textit{Initialize\_population\_with\_optimization} ($P$, $\mathbf{X}$, $p$); 
		\FOR{$i$ from $1$ to $\mathcal{T}$}
		\STATE \textit{Compute\_perplexity} (LLM, $\text{P}_{i-1}$, $\mathbf{X}$);
		\STATE Top-k $\gets$ \textit{Update\_Topk} (Top-k, $\text{P}_{i-1}$);
		\STATE $\text{P}_{mutation}$ $\gets$ \textit{Mutation\_with\_mono\_constraint} (Top-k, $N_1$, $p$); 
		\STATE $\text{P}_{crossover}$ $\gets$ \textit{Crossover\_with\_mono\_constraint} (Top-k, $N_2$);
		\STATE $\text{P}_i$ $\gets$ $\text{P}_{mutation}$ $\cup$ $\text{P}_{crossover}$ $\cup$ Top-k;
		\ENDFOR
		\STATE Output the candidate from Top-k with minimal perplexity;
	\end{algorithmic}
\end{algorithm}

\textbf{8$\times$ extension without fine-tuning}. By employing an evolutionary search, we successfully determine non-uniform RoPE rescale factors that retain essential dimensions and positional information, thereby reducing the information loss typically caused by interpolation. As illustrated in Fig.\ref{fig:non-ft}, this strategy enables LLaMA2's context window to be expanded from 4k to 32k tokens without the need for fine-tuning. In comparison, current approaches like PI, as well as non-uniform NTK and YaRN, experience significant increases in perplexity after only a 2$\times$ context extension.

\subsection{Extending LLM Context Window to 2048K}
\label{sec:extendto2048k}

\textbf{Progressive extension to 2048k}. Here, we present an approach to expand the context window of pre-trained LLMs from the standard 4k to beyond 2048k tokens. As shown previously, applying non-uniform positional interpolation facilitates an 8$\times$ increase without the need for fine-tuning. However, when attempting greater extensions—such as scaling up to 512$\times$—fine-tuning becomes indispensable. One possible strategy involves identifying optimal RoPE rescale factors for the intended 2048k context length, followed by subsequent fine-tuning. Nevertheless, this method encounters difficulties due to the significant computational demands of such large-scale training. Furthermore, our experience suggests that effectively fine-tuning LLMs for such substantial context extensions is quite problematic (refer to Appendix).

Fortunately, LongRoPE demonstrates efficacy when applied to both the base and the adapted extended LLM. To this end, we propose a progressive and efficient strategy that reaches the desired 2048k target using merely 1k fine-tuning iterations, all confined to a maximum training length of 256k.

$\Diamond$ \textit{Investigating LongRoPE search for augmenting pre-trained LLM to 256k}. Using LLaMA2 as the illustrative case, we explore different configurations aimed at achieving context window capacities of 128k and 256k. This phase involves scaling up by a factor of 32$\times$ and 64$\times$, respectively.

$\Diamond$ \textit{Fine-tuning to 256k}. Subsequently, the pre-trained LLM undergoes fine-tuning in order to extend its context window to 256k tokens. The procedure involves an initial 400 steps of fine-tuning LLaMA2 utilizing RoPE rescaled factors set to 128k. Afterward, the RoPE rescaled factors are updated to 256k on the resulting checkpoint, followed by a further 600 steps of fine-tuning. Compared to fine-tuning directly to 256k, this staged strategy demonstrates greater efficiency.

$\Diamond$ \textit{Increasing the context window of the fine-tuned extended LLM to 2048k via LongRoPE search}. In the final stage, an additional search procedure is applied to the LLM already fine-tuned for 256k sequence lengths. This process successfully scales the context window to an impressive size of 2048k, accomplished without any subsequent fine-tuning. The total enlargement factor reaches $512\times$.

\textbf{Original context window degradation}.  Expanding the context window to a substantial 2048k length results in reduced performance when operating within the initial context range. This phenomenon, which has been documented as a drawback of positional interpolation~\cite{pi}, stems from the compression of positional embeddings for the original context window into a significantly smaller subspace. Such compression adversely influences the effectiveness of the language model. Specifically, when a 512$\times$ scaling is applied, the position representations corresponding to the original 4k window are densely concentrated, exacerbating the issue.

To address this issue, we conduct an additional evolutionary search targeting the expanded LLM, specifically to fine-tune the RoPE rescale parameters for brief context lengths such as 4k and 8k tokens. The search space for the maximum permissible $\lambda$ is constrained, reflecting the decreased necessity for positional interpolation when handling shorter sequences. In the inference phase, the LLM adaptively modifies the relevant RoPE rescale factors in response to the current context length.

 \section{Experiments}

\subsection{Setup}
\textbf{Evaluation Tasks and models}. In this study, LongRoPE is integrated with LLaMA2-7B and Mistral-7B, and its efficacy is assessed across three dimensions: (1) the perplexity scores of large language models with extended context lengths when processing lengthy documents; (2) a Passkey retrieval task designed to test how effectively the model extracts a straightforward passkey from abundant irrelevant information; and (3) traditional LLM benchmark tests conducted using a restricted context window of 4096 tokens.

\textbf{Fine-tuning}. For LLaMA2, we adopt a linear learning rate schedule starting at 2e-5 and utilize a global batch size of 32. The initial fine-tuning phase spans 400 steps using the Redpajama~\cite{redpajama} dataset, organized into 128k-length segments that are enclosed with BOS and EOS tokens. Building upon this pretrained checkpoint, we extend training by an additional 600 steps to expand to a 256k context window. Training at the 128k context is performed on 8 A100 GPUs via the distributed training platform~\cite{cube}, whereas the 256k context window requires scaling up to 16 A100 GPUs. For Mistral, we employ a constant learning rate of 1e-6 and set the global batch size to 64. For both the 128k and 256k variants, our methodology adheres to the YaRN configuration~\cite{yarn}, entailing 400 training steps on Together Computer's Long-Data Collections~\cite{mistraldata} with sequences of length 16k, leveraging 4 A100 GPUs in all cases.

\textbf{Search}. When the target window size is at or below 256k, the parameters are set as follows: $P$=64, $N_1$=$N_2$=16, $p$=0.3, and $\mathcal{T}$=40. In each search step, the top 32 candidates are chosen for mutation and crossover. To evaluate perplexity, five random validation samples are drawn from PG19, each meeting a minimum threshold of the intended context length. For window sizes exceeding 512k, the population size and the numbers for mutation and crossover are reduced by 50\%. In this regime, perplexity assessment uses three random validation samples from the Pile-Books3~\cite{pile} dataset.

\textbf{Baselines}. To achieve a 2048k context window, models with initial context windows of 128k and 256k were fine-tuned, resulting in LongRoPE-2048k (ft=128k) for LLaMA2 and LongRoPE-2048k (ft=256k) for Mistral. These four variants are evaluated against leading approaches for context window extension in open-source LLMs, specifically those fine-tuned post-positional interpolation using methods such as PI, NTK, and YaRN. Among the baseline models considered are Together-32k~\cite{together}, Code LLaMA~\cite{codellama}, LongLoRA-full-FT-100k~\cite{longlora}, YaRN-LLaMA, and YaRN-Mistral~\cite{yarn}.

\subsection{Main Results}

\label{sec:mainresult}
\textbf{Language modeling over sequences up to 256k}. Our initial analysis contrasts our method with leading long-context language models, focusing on an evaluation length of 256k. To showcase the robustness of our approach, we employ two benchmark datasets: the test splits from Proof-pile~\cite{pg19} and PG19~\cite{pile}. We assess model performance in terms of perplexity across a range of context lengths, employing a sliding window of size 256. For PG19, the evaluation is conducted over its complete test split, comprising 100 documents. In line with the procedure established by YaRN~\cite{yarn}, we sample 10 sequences, each exceeding 128k tokens in length, from Proof-pile for our experiments.

Table~\ref{tbl:proof-pile} and Table~\ref{tbl:pg19} present a perplexity comparison of LLaMA2 and Mistral, both adapted using various interpolation strategies, on the Proof-pile and PG19 datasets, respectively. Two main findings emerge: \textbf{(1)} the adapted models demonstrate a clear reduction in perplexity as evaluation length increases from 4k to 256k, indicating improved ability to utilize extended context. \textbf{(2)} Notably, when evaluated at a 16$\times$ extended context window—a scenario that commonly hinders retention of performance for shorter sequences—our LongRoPE-2048k variants still surpass leading baseline models within the 256k context window.

\textbf{Language modeling for extended sequences exceeding 2000k}. To assess performance on exceptionally lengthy texts, we utilize the Books3~\cite{pile} dataset. For efficient evaluation, 20 books are randomly chosen, each with a length greater than 2048k, and a 256k sliding window is applied.

Referencing Table~\ref{tbl:books3}, LongRoPE is able to expand the context length of both LLaMA2-7B and Mistral-7B to 2048k tokens, while maintaining or surpassing baseline perplexity values within the 8k-128k range. Significant discrepancies are evident when comparing LLaMA2 and Mistral at the 2048k context size. Mistral generally exceeds baseline performance in shorter contexts, yet its perplexity rises above 7 beyond the 256k mark. LLaMA2 displays the anticipated trend, with perplexity decreasing as the context grows longer, albeit with slight increases at 1024k and 2048k. Notably, for LLaMA2, LongRoPE-2048k yields improved results when fine-tuned at 256k, as opposed to 128k, owing to the lower secondary extension ratio (i.e., 8$\times$ versus 16$\times$). Conversely, Mistral achieves superior outcomes when fine-tuning at a window size of 128k. This divergence stems largely from following YaRN's configuration during Mistral's 128k and 256k fine-tuning, which employs a training length of 16k and consequently impacts the model's capacity for additional context extension after fine-tuning.

\textbf{Passkey retrieval}. Next, we examine how the context window size impacts performance in generative tasks. We adopt the synthetic passkey retrieval benchmark introduced by~\cite{passkey}, where the objective is to identify a concealed passkey—a random five-digit sequence—embedded somewhere within an extensive document. The structure of the prompt utilized can be found in the appendix. For evaluation, we conduct 10 trials of the passkey retrieval task, each time positioning the passkey at a randomly selected location, drawn uniformly throughout the evaluation context window.

Fig.~\ref{fig:passkey} presents a comparison of retrieval accuracy against baseline models. While prior approaches exhibit a sharp decline in accuracy to zero past the 128k token length, our LongRoPE-LLaMA2-2048k (ft=256k) consistently sustains high retrieval performance ($\ge$90\%) across sequence lengths ranging from 4k up to 2048k, despite the inherently difficult nature of retrieving a passkey from sequences on the scale of millions of tokens. For LongRoPE-Mistral-2048k (ft=128k), accuracy is perfectly retained up to 1800k tokens, only dropping to 60\% at 2048k. This decrease is in accordance with Table~\ref{tbl:books3}, which demonstrates a modest rise in perplexity at 2048k tokens.

\textbf{Evaluation on established benchmarks within the native context window}. The performance of LongRoPE-2048k models is assessed using the Hugging Face Open LLM Leaderboard~\cite{huggingfaceboard}, considering both zero-shot and few-shot scenarios. Specifically, the experiments utilize 25-shot settings for ARC-Challenge~\cite{allenai:arc}, 10-shot configurations for HellaSwag~\cite{zellers2019hellaswag}, 5-shot instances for MMLU~\cite{mmlu}, and a 0-shot setting for TruthfulQA~\cite{lin2021truthfulqa}.

Table~\ref{tbl:commonsense} illustrates that our models deliver similar outcomes to those obtained on the original benchmark, which targets shorter context lengths, and surpass the initial Mistral model on TruthfulQA with a gain of +0.5\%. LongRoPE-LLaMA2-2048k, tuned at 256k, experiences a modest drop in performance, yet stays largely consistent across the majority of evaluations.

\subsection{Ablation Results}

\textbf{Evaluating the effectiveness of secondary positional interpolation}. Within the framework of our progressive extension methodology, the search algorithm facilitates a second phase of non-uniform positional interpolation applied to the fine-tuned, extended LLMs. To assess its impact, we conduct experimental tests using the fine-tuned LLaMA2-256k model, expanding its context window to 512k, 1024k, and 2048k via PI and YaRN. According to the results in Table~\ref{tbl:secondsearch}, our approach of non-uniform positional interpolation preserves stable perplexity across extensions. In contrast, the perplexity scores with PI and YaRN escalate rapidly as the extension ratio increases.

\textbf{Recovery performance at reduced context lengths.} In order to address performance degradation when operating at shorter context windows, we apply our search algorithm to modify the RoPE scaling for LongRoPE-2048k. Concretely, we restrict the search space by lowering the maximum scale factors, thereby minimizing interpolation for 4k and 8k contexts. 
Table~\ref{tbl:shortperformance} presents perplexity scores for LongRoPE-LLaMA2-2048k evaluated on Proof-pile at 4k and 8k contexts, alongside the mean LLM benchmark accuracy. The outcomes indicate that these adjustments yield considerable improvements in model performance at reduced context lengths.

\textbf{Examination of the two distinct non-uniformities}. In our final ablation study, we investigate how each non-uniformity individually affects model outcomes. We design two experimental setups: (i) expanding LLaMA2-7B to shorter sequence lengths of 16k and 32k utilizing varying approaches—PI, tuning solely the RoPE dimension, and adjusting both forms of non-uniformity; (ii) scaling our fine-tuned LLaMA2 model with 256k-length inputs up to 2048k in sequence, using identical procedures. Perplexity is computed without any subsequent fine-tuning. According to Table~\ref{tbl:searchdimension}, integrating non-uniformity in the RoPE dimension substantially lowers perplexity compared to linear interpolation with PI. Additionally, introducing non-uniformity in the token position leads to clear gains for 16k and 32k sequence lengths, but yields less noticeable improvements at 2048k, likely due to the challenges associated with ultra-long contexts. Simple retention of initial tokens, absent interpolation, offers limited value in these extreme cases; we propose exploring this further in future research.

\section{Related Works}

Besides position interpolation methods, this section also reviews relevant studies employing alternative strategies.

\textbf{Retrieval-based} techniques incorporate an external memory component to store extended historical context and deploy retrieval modules to access pertinent documents during inference~\cite{longllama,wang2023augmenting,borgeaud2022improving}. These solutions often require substantial changes to the underlying LLM architectures. In contrast, our approach introduces minimal modifications, confined primarily to positional embeddings, making it notably lighter. Furthermore, our method extends to a wider array of long context tasks beyond just retrieval, including applications such as long document summarization and few-shot learning.

\textbf{Attention-driven expansion of context windows}. In addition to approaches relying on positional embedding interpolation, several studies extend the input context window for LLMs by directly modifying attention mechanisms~\cite{han2023lminfinite, streamingllm,parallelwindow}. The primary strategy involves addressing the growth in attention computation from novel positions through innovative attention mask designs. These attention-focused techniques are designed to work in tandem with positional interpolation methods, offering complementary benefits.

\textbf{Fine-tuning based approaches} investigate strategies to adapt pre-trained LLMs to handle longer context windows through modified position embeddings. Studies such as Code LLaMA~\cite{codellama}, LLaMA2 Long~\cite{xiong2023effective}, and ScaledRoPE~\cite{liu2023scaling} employ significantly increased base values for RoPE, followed by fine-tuning at the desired context length. In contrast, our approach provides adaptability across a range of target lengths and can scale beyond lengths of 2M. Recently, due to the high computational burden of fine-tuning at very long sequence lengths (i.e., beyond 128k), methods like LongLoRA~\cite{longlora} and PoSE~\cite{zhu2023pose} have been introduced to reduce GPU requirements. Our approach is complementary and independent of these resource-efficient fine-tuning techniques.

\section{Conclusion}

This paper introduces LongRoPE, a technique that substantially increases the context length of LLMs up to an impressive 2048k tokens, while preserving performance within their original shorter window. Our approach leverages two distinct types of non-uniformities found in RoPE positional embeddings, guided by an efficient evolutionary search process. This method yields dual advantages: it achieves strong initialization for subsequent fine-tuning and facilitates up to an 8$\times$ expansion of the context window even without fine-tuning. Furthermore, we implement a staged extension strategy, employing fine-tuned LLMs with a 256k context window to scale up to 2048k tokens, all without additional fine-tuning. Comprehensive experimental evaluation demonstrates LongRoPE’s efficacy. We anticipate LongRoPE-2048k models will unlock a range of novel applications requiring extended context and motivate continued advances in the field.

\section*{Broader Impacts} The research outlined in this paper seeks to further progress in the domain of Machine Learning. Although our contributions may carry various implications for society, we do not identify any particular impact that warrants explicit attention in this section.

	\balance

\end{document}